\backmatter

\chapter{Afterword: A Final Thought}

\marksourcepage{806}
\noindent A Final Thought\par
Because the use of a theoretical pricing model requires a trader to make so many different decisions with respect to both the inputs into the model and the reliability of the assumptions on which the model is based, a new option trader may feel that making the right decisions is either an impossible task or simply a matter of luck. It is true that a trader using a model will almost certainly be wrong about at least some of the inputs into the model, and luck undoubtedly does play a role in the short run. But in the long run traders who are willing to put in the effort required to understand how a model works, including its strengths and weaknesses, always seem to come out ahead. Experienced traders know that under most conditions, using a model, with all its problems, is still the best way to evaluate options and manage risk.\par
Regardless of whether a model is simple or complex, a trader who uses\par
a model needs to have faith in the model. Otherwise, why use a model at all? Indeed, for traders who are not mathematically proficient, using a model is often a leap of faith. But having faith in a model does not mean having blind, unquestioning faith. If a model returns values that are clearly inconsistent with common sense, or if market conditions are changing so quickly that it is impractical to use the model in its current form, a trader may have to decide whether to adjust the model, if that is possible, or simply to stop using the model. Although we have emphasized the importance of models, trading is both an art and a science. Experienced traders know that there are times when it is perhaps best to put the model aside and make decisions based on other intangible assets, whether intuition, ``market feel,'' or experience. A trader who slavishly uses a model to make every trading decision is heading for disaster. Only a trader who fully understands what a model can and cannot do will be able to make the model his servant rather than his master.\par

\chapter{Glossary of Option Terminology}

\marksourcepage{807}
\noindent A\par
This glossary includes option-related terms as they are most commonly\par
used. However, the reader should be aware that option terminology is not uniform. Traders may sometimes refer to different strategies or option characteristics with the same term. They may sometimes refer to the same strategy or characteristic with different terms.\par
All or None (AON) An order that must be filled in its entirety or not at all.\par
American Option An option that can be exercised at any time prior to\par
expiration.\par
Arbitrage The purchase and sale of the same or closely related products in different markets to take advantage of a price disparity between the two markets.\par
Asian Option See Average Price Option.\par
Assignment The process by which the seller of an option is notified of the buyer's intention to exercise. The seller is required to take a short position in the underlying position in the case of a call or a long position in the case of a put.\par
At the Forward An option whose exercise price is equal to the forward\par
price of the underlying contract. Sometimes referred to as At-the-Money Forward.\par

\marksourcepage{808}
At the Money An option whose exercise price is equal to the current price of the underlying contract. On listed option exchanges, the term is more commonly used to refer to the option whose exercise price is closest to the current price of the underlying contract.\par
Automatic Exercise The exercise by the clearinghouse of an in-the-money option at expiration unless the holder of the option submits specific instructions to the contrary.\par
Average Price Option An option whose value at expiration is determined by the average price of the underlying instrument over some period of time. Also known as an Asian Option.\par
Backspread A spread, usually  delta  neutral, where more options are\par
purchased than sold, where all options are the same type and expire at the same time.\par
Backward A futures market where the long-term delivery months trade at a discount to the short-term delivery months.\par
Barrier Option A type of exotic option that will either become effective or cease to exist if the underlying instrument trades at or beyond some predetermined price prior to expiration.\par
Bear Spread Any spread that will theoretically increase in value with a decline in the price of the underlying contract.\par
Bermuda Option An option that can be exercised prior to expiration, but only during a predetermined period or window. Also known as a Mid- Atlantic Option.\par
Binary Option An option that, if in the money at expiration, makes one predetermined payout. Also known as a Digital Option.\par
Box A long call and short put at one exercise price, together with a short call and long put at a different exercise price. All options must have the same underlying contract and expire at the same time.\par

\marksourcepage{809}
Bull Spread Any spread that will theoretically increase in value with a rise in the price of the underlying contract.\par
Butterfly The sale (purchase) of two options with the same exercise price, together with the purchase (sale) of one option with a lower exercise price and one option with a higher exercise price. All options must be of the same type, have the same underlying contract, and expire at the same time, and there must be an equal increment between exercise prices.\par
Buy/Write The purchase of an underlying contract together with the sale of a call option on that contract.\par
Cabinet Bid An option price that is smaller than the normally allowable minimum price. On some exchanges, a cabinet bid is permissible between traders desiring to close out positions in very far out-of-the-money options.\par
Calendar Spread The purchase (sale) of one option expiring on one date and the sale (purchase) of another option expiring on a different date. Typically, both options are of the same type, have the same exercise price, and have the same underlying stock or commodity. Also known as a Time Spread or Horizontal Spread.\par
Call Option A contract between a buyer and a seller whereby the buyer\par
acquires the right, but not the obligation, to purchase a specified underlying contract at a fixed price on or before a specified date. The seller of the call option assumes the obligation of delivering the underlying contract should the buyer wish to exercise his option.\par
Cap A contract between a borrower and a lender of floating-rate funds\par
whereby the borrower is assured of paying no more than some maximum\par
interest rate for borrowed funds. This is analogous to a call option where the underlying instrument is an interest rate on borrowed funds.\par
Charm The sensitivity of an option's delta to the passage of time.\par
Chooser Option A straddle where  the owner must decide by some\par
predetermined date whether to keep either the call or the put.\par

\marksourcepage{810}
Christmas Tree A spread involving three exercise prices. One or more calls (puts) are purchased at the lowest (highest) exercise price, and one or more calls (puts) are sold at each of the higher (lower) exercise prices. All options must expire at the same time, be of the same type, and have the same underlying contract. Also known as a Ladder.\par
Class All options of the same type with the same expiration date and same underlying instrument.\par
Clearinghouse The organization that guarantees the integrity of all trades made on an exchange.\par
Clearing Member A member firm of an exchange that is authorized by the clearinghouse to process trades for its customers and that guarantees, through the collection of margin and variation, the integrity of its customers' trades.\par
Collar A long (short) underlying position that is hedged with both a long (short) out-of-the-money put and a short (long) out-of-the-money call. All options must expire at the same time. Also known as a Cylinder, Fence, or Range Forward.\par
Color The sensitivity of an option's gamma to the passage of time.\par
Combination (Combo) A two-sided option spread that does not fall into\par
any well-defined category of spreads. Most commonly, it refers to a long call and short put or short call and long put, which together make up a synthetic position in the underlying contract.\par
Compound Option An option to purchase an option.\par
Condor The sale (purchase) of two options with different exercise prices, together with the purchase (sale) of one option with a lower exercise price and one option with a higher exercise price. All options must be of the same type, have the same underlying contract, and expire at the same time, and there must be an equal increment between exercise prices.\par

\marksourcepage{811}
Contango A futures market where the long-term delivery months trade at a premium to the short-term delivery months.\par
Contingency Order An  order  that becomes  effective  only on  the\par
fulfillment of some predetermined condition(s) in the marketplace.\par
Conversion A long underlying position together with a synthetic short\par
underlying position. The synthetic position consists of a short call and long put, where both options have the same exercise price and expire at the same time. Sometimes referred to as a Forward Conversion.\par
Covered Write The sale of a call (put) option against an existing long (short) position in the underlying contract.\par
Cylinder See Collar.\par
Delta (Delta) The sensitivity of an option's theoretical value to a change in the price of the underlying contract. Also known as the Hedge Ratio.\par
Delta Neutral A position where the sum total of all the deltas add up to approximately 0. Under current market conditions, the position has no preference as to the direction of movement in the underlying market.\par
Diagonal Roll See Time Box.\par
Diagonal Spread A long option at one exercise price and expiration date, together with a short option at a different exercise price and expiration date. All options must be the same type. This is the same as a calendar spread using different exercise prices.\par
Digital Option See Binary Option.\par
Dragonfly A long (short) straddle, together with two short (long) strangles at the same exercise price, where all options expire at the same time and have the same underlying contract. The exercise price of the straddle will usually fall as close as possible to the midpoint between the exercise prices of the strangles.\par

\marksourcepage{812}
Dynamic Hedging A  process  in  which  the  underlying  contract  is\par
periodically bought or sold in order to maintain a desired position in a market. Dynamic hedging is most often used to maintain a delta-neutral option position.\par
Efficiency A number that represents the relative risk and reward of a\par
potential option strategy. The risk and reward are typically represented by the total gamma, theta, and vega of the strategy. The efficiency is generated by dividing one sensitivity by another.\par
Elasticity The percent change in an option's value for a given percent change in the value of the underlying instrument. Sometimes referred to as an option's Leverage Value. The elasticity is sometimes denoted by the Greek letter Lambda (Lambda).\par
Eurocurrency Currency deposited in a bank outside the currency's home\par
country.\par
Eurocurrency rate The interest rate paid on currency deposited in a bank outside the currency's home country.\par
European Option An option that may only be exercised at expiration.\par
Exchange Option An option to exchange one asset for another asset.\par
Ex-Dividend The first day on which a dividend-paying stock is trading\par
without the right to receive the dividend.\par
Exercise The process by which the holder of an option notifies the seller of his intention to take a long position in the underlying contract in the case of a call or a short position in the underlying contract in the case of a put.\par
Exercise Price The price at which the underlying contract will be delivered in the event an option is exercised. Also known as the Strike Price.\par
Expiration (Expiry) The date and time after which an option may no\par
longer be exercised.\par

\marksourcepage{813}
Exotic Option An option with nonstandard contract specifications. Sometimes referred to as a Second-Generation Option. Exotic options are usually traded in the over-the-counter (off -exchange) market.\par
Extrinsic Value See Time Value.\par
Fair Value See Theoretical Value.\par
Fence See Collar.\par
Fill or Kill (FOK) An order that will automatically be canceled unless it can be executed immediately and in its entirety.\par
Flex Option An exchange-traded option where the buyer and seller are\par
permitted to negotiate the exact terms of the option contract. Typically, this includes the exercise price, the expiration date, and the exercise style (either European or American).\par
Floor A contract between a borrower and a lender of floating-rate funds whereby the lender is assured of receiving no less than some minimum interest rate for loaned funds. This is analogous to a put option where the underlying instrument is an interest rate on loaned funds.\par
Forward Contract An agreement between a buyer and a seller to exchange money for goods at some later date. At maturity, the buyer is obligated to take delivery and the seller is obligated to make delivery.\par
Forward Conversion See Conversion.\par
Forward Price The price that the buyer of a forward contract agrees to pay at maturity of the contract.\par
Forward Start Option An option that becomes effective only on some\par
future predetermined date.\par
Front Spread A spread, usually delta neutral, where more options are sold than purchased, where all options are the same type, and all expire at the same time.\par

\marksourcepage{814}
Fugit Assuming that all market conditions remain unchanged, the expected amount of time remaining to optimal early exercise of an American option.\par
Futures Contract An exchange-traded forward contract.\par
Futures-Type Settlement A settlement procedure used by commodity\par
exchanges whereby an initial margin deposit is made but under which no immediate cash payment is made by the buyer to the seller. Cash settlement takes place at the end of each trading day based on the difference between the original trade price or the previous day's settlement price and the current day's settlement price.\par
Gamma (?) The sensitivity of an option's delta to a change in the price of the underlying contract.\par
Good `til Canceled (GTC) An order that remains active until it can either be executed or is canceled by the customer.\par
Guts A strangle where both the call and the put are in the money.\par
Haircut On a securities exchange, money that a professional trader is required to keep in his account in order to cover the risk of his position. Haircut requirements are normally determined by the regulatory authority under which the exchange operates.\par
Hedge Ratio See Delta.\par
Hedger A trader who enters the market with the specific intent of protecting an existing position in an underlying contract.\par
Horizontal Spread See Calendar Spread.\par
Immediate or Cancel (IOC) An order that will automatically be canceled if it cannot be executed immediately. An IOC order need not be filled in its entirety.\par
Implied Volatility Assuming that all other inputs are known, the volatility that would have to be input into a theoretical pricing model to yield a\par

\marksourcepage{815}
theoretical value that is identical to the price of the option in the marketplace.\par
In-Option A barrier option that becomes effective only if the underlying instrument trades at or through some predetermined price prior to expiration. Also known as a Knock-In Option.\par
In-Price The price at which the underlying instrument must trade before an in-option becomes effective.\par
In the Money An option that has intrinsic value greater than 0. A call option is in the money if its exercise price is lower than the current price of the underlying contract. A put option is in the money if its exercise price is higher than the current price of the underlying contract. An option may also be In the Money Forward if it has intrinsic value greater than 0 when compared with the forward price of the underlying contract.\par
Index Arbitrage A strategy which attempts to profit from the relative mispricing of options, futures contracts, or the component stocks which make up a stock index.\par
Intermarket Spread A spread consisting of opposing market positions in two or more different underlying securities or commodities.\par
Intrinsic Value For an in-the-money option, the difference between the exercise price and the underlying price. Out-of-the-money options have no intrinsic value. An option whose price is equal to its intrinsic value is said to be trading at Parity.\par
Iron Butterfly A long (short) straddle, together with a short (long) strangle, where all options expire at the same time and have the same underlying contract. The exercise price of the straddle is located at the midpoint between the exercise prices of the strangle.\par
Iron Condor A long (short) strangle with narrower exercise prices, together with a short (long) strangle with wider exercise prices, where all options expire at the same time and have the same underlying contract. The\par

\marksourcepage{816}
narrower strangle is centered between the exercise prices of the wider strangle.\par
Jelly Roll See Roll.\par
Kappa (K) See Vega. The Greek letter kappa is sometimes used to denote an option's exercise price.\par
Knock-In Option See In-Option.\par
Knock-Out Option See Out-Option.\par
Ladder See Christmas Tree. Alternatively, a type of exotic option whose minimum value increases as the underlying contract goes through a series of predetermined prices, or rungs, over the life of the option.\par
Lambda (Lambda) See Elasticity.\par
LEAP (Long-Term Equity Anticipation Security) A long-term (usually\par
more than one year) exchange-traded equity option.\par
Leg One side of a spread position.\par
Leverage Value See Elasticity.\par
Limit The maximum allowable price movement over some time period for\par
an exchange-traded contract.\par
Limit Order An order to be executed at a specified price or better.\par
Local An independent trader on a commodity exchange. Locals perform\par
functions similar to market makers on stock and stock option exchanges.\par
Locked Market An exchange-traded market where trading has been halted\par
because prices have reached the limit permitted by the exchange.\par
Long A position resulting from the purchase of a contract. The term is also used to describe a position that will theoretically increase (decrease) in\par

\marksourcepage{817}
value should the price of the underlying contract rise (fall). Note that a long (short) put position is a short (long) market position.\par
Long Premium A position that will theoretically increase in value should the underlying contract make a large or swift move in either direction. The position will theoretically decrease in value should the underlying market fail to move or move very slowly. The term may also refer to a position that will increase in value should implied volatility rise.\par
Long Ratio Spread A spread where more options are purchased than sold.\par
Lookback Option An exotic option whose exercise price will be equal to either the lowest price of the underlying instrument in the case of a call or the highest price of the underlying instrument in the case of a put over the life of the option. A lookback option can also have a fixed strike, in which case its value at expiration will be determined by the maximum underlying price in the case of a call or the minimum underlying price in the case of a put over the life of the option.\par
Margin Money deposited by a trader with the clearing house to ensure the integrity of his trades.\par
Market-if-Touched (MIT) A contingency order that becomes a market\par
order if the contract trades at or beyond a specified price.\par
Market maker An independent trader or trading firm, usually appointed by an exchange, that is prepared to both buy and sell contracts in a designated market. A market maker is required to quote both a bid and offer price in his designated contract.\par
Market-on-Close (MOC) An order to be executed at the market price at\par
the close of that day's trading.\par
Market Order An order to be executed immediately at the current market price.\par
Mark-to-market A method of valuing a position based on the current\par
market price of all contracts which make up the position.\par

\marksourcepage{818}
Married Put A long (short) put together with a long (short) underlying position.\par
Mid-Atlantic Option See Bermuda Option.\par
Midcurve Option In futures option markets, a short-term option on a longterm futures contract. Midcurve options are most common in euro-currency futures markets, such as Eurodollars and Euribor.\par
Naked A long (short) market position with no offsetting short (long) market position.\par
Neutral Spread A spread that is neutral with respect to some risk measure, most commonly the delta. A spread may also be lot neutral, where the total number of long and short contracts of the same type are equal.\par
Not Held An order submitted to a broker but over which the broker has\par
discretion as to when and how the order is executed.\par
Omega (Omega) The Greek letter sometimes used to denote an option's elasticity. An alternative to lambda (Lambda).\par
One-Cancels-the-Other (OCO) Two orders submitted simultaneously, either of which may be executed. If one order is executed, the other is automatically canceled.\par
Order Book  Official (OBO) An  exchange  official  responsible  for\par
executing market or limit orders for public customers.\par
Out of the Money An option that currently has no intrinsic value. A call is out of the money if its exercise price is more than the current price of the underlying contract. A put is out of the money if its exercise price is less than the current price of the underlying contract. An option may also be Out of the Money Forward if it has no intrinsic value when compared with the forward price of the underlying contract.\par
Out-Option A type of barrier option that is deemed to have expired if the underlying instrument trades at some predetermined price prior to\par

\marksourcepage{819}
expiration. Also known as a Knock-Out Option.\par
Out-Price The price at which the underlying instrument must trade before an out-option is deemed to have expired.\par
Out-Trade A trade that cannot be processed by the clearinghouse due to conflicting information reported by the two parties to the trade.\par
Overwrite The sale of an option against an existing position in the underlying contract.\par
Parity See Intrinsic Value.\par
Phi (?) For foreign-currency options, the sensitivity of the option's value to a change in the foreign interest rate. Sometimes referred to as Rho2.\par
Pin Risk The risk to the seller of an option that at expiration will be exactly at the money. The seller will not know whether the option will be exercised.\par
Portfolio Insurance A process in which the quantity of holdings in an\par
underlying instrument is periodically adjusted to replicate the characteristics of an option on the underlying instrument. This is similar to the delta-neutral dynamic hedging process used to capture the value of a mispriced option.\par
Position The sum total of a trader's open contracts in a particular underlying market.\par
Position Limit For an individual trader or firm, the maximum number of open contracts in the same underlying market permitted by an exchange or clearinghouse.\par
Premium The price of an option.\par
Program Trading An arbitrage strategy involving the purchase or sale of a stock index futures contract against an opposing position in the component stocks that make up the index.\par

\marksourcepage{820}
Put Option A contract between a buyer and a seller whereby the buyer\par
acquires the right but not the obligation to sell a specified underlying contract at a fixed price on or before a specified date. The seller of the put option assumes the obligation of taking delivery of the underlying contract should the buyer wish to exercise his option.\par
Range Forward See Collar.\par
Ratchet Option A type  of  exotic  option whose minimum  value  is\par
determined by the underlying price  at a series of predetermined time\par
intervals over the life of the option.\par
Ratio Spread Any spread where the number of long market contracts (long underlying, long call, or short put) and short market contracts (short underlying, short call, or long put) are unequal.\par
Ratio Write The sale of multiple options against an existing position in an underlying contract.\par
Reversal See Reverse Conversion.\par
Reverse Conversion A short underlying position together with a synthetic long underlying position. The synthetic position consists of a long call and short put, where both options have the same exercise price and expire at the same time. Also known as a Reversal.\par
Rho (P) The sensitivity of an option's theoretical value to a change in interest rates.\par
Risk Reversal A long (short) underlying position together with a long\par
(short) out-of-the-money put and a short (long) out-of-the-money call. Both options must expire at the same time. Also known as a Split-Strike Conversion. The position is equivalent to a Collar.\par
Roll A long call and short put with one expiration date, together with a short call and long put with a different expiration date. All four options must have the same exercise price and the same underlying stock or commodity. In slang terms, sometimes referred to as a Jelly Roll.\par

\marksourcepage{821}
Scalper A floor trader on an exchange who hopes to profit by continually buying at the bid price and selling at the offer price in a specific market. Scalpers usually try to close out all positions at the end of each trading day.\par
Second-Generation Option See Exotic Option.\par
Serial Option On  futures  exchanges, an  option  expiration  with no\par
corresponding futures expiration. The underlying contract for a serial option is the nearest futures contract beyond the option expiration.\par
Series All options with the same underlying contract, same exercise price, and same expiration date.\par
Short A position resulting from the sale of a contract. The term is also used to describe a position that will theoretically increase (decrease) in value should the price of the underlying contract fall (rise). Note that a short (long) put position is a long (short) market position.\par
Short Premium A position that will theoretically increase in value should the underlying contract fail to move or move very slowly. The position will theoretically decrease in value should the underlying market make a large or swift move in either direction. The term may also refer to a position that will increase in value should implied volatility fall.\par
Short Ratio Spread A spread where more options are sold than purchased.\par
Short Squeeze A situation in the stock option market, usually resulting from a partial tender offer, where no stock can be borrowed to maintain a short stock position. If assigned on a short call position, a trader may be forced to exercise a call early to fulfill his delivery obligations, even though the call still has some time value remaining.\par
Sigma (sigma) The commonly used notation for standard deviation. Because volatility is usually expressed as a standard deviation, the same notation is often used to denote volatility.\par
Specialist A market maker given exclusive rights by an exchange to make a market in a specified contract or group of contracts. A specialist may buy or\par

\marksourcepage{822}
sell for his own account or act as a broker for others. In return, a specialist is required to maintain a fair and orderly market.\par
Speculator A trader who hopes to profit from a specific directional move in an underlying contract.\par
Speed The sensitivity of an option's gamma to a change in the underlying price.\par
Spread A long market position and an offsetting short market position\par
usually, but not always, in contracts with the same underlying market.\par
Split-Strike Conversion See Risk Reversal.\par
Stock-Type Settlement A settlement procedure in which the purchase of a contract requires full and immediate payment by the buyer to the seller. All profits or losses from the trade are unrealized until the position is liquidated.\par
Stop-Limit Order A contingency order that becomes a limit order if the contract trades at a specified price.\par
Stop (Loss) Order A contingency order that becomes a market order if the contract trades at a specified price.\par
Straddle A long (short) call and a long (short) put where both options have the same underlying contract, the same expiration date, and the same exercise price.\par
Strangle A long (short) call and a long (short) put where both options have the same underlying contract, the same expiration date, but different exercise prices.\par
Strap An archaic term for a position consisting of two long (short) calls and one long (short) put where all options have the same underlying contract, the same expiration date, and the same exercise price.\par
Strike Price (Strike) See Exercise Price.\par

\marksourcepage{823}
Strip An archaic term for a position consisting of one long (short) call and two long (short) puts where all options have the same underlying contract, the same expiration date, and the same exercise price. Alternatively, in Eurocurrency markets, a series of futures or futures options designed to replicate the characteristics of a long-term interest-rate position.\par
Swap An agreement to exchange cash flows. Most commonly, a swap\par
involves exchanging variable-interest-rate payments for fixed-interest-rate payments.\par
Swaption An option to enter into a swap agreement.\par
Synthetic A combination of contracts that together have approximately the same characteristics as some other contract.\par
Synthetic Call A long (short) underlying position together with a long (short) put.\par
Synthetic Put A short (long) underlying position together with a long\par
(short) call.\par
Synthetic Underlying A long (short) call and short (long) put where both options have the same underlying contract, the same expiration date, and the same exercise price.\par
Tau (tau) The commonly used notation for the amount of time remaining to expiration. Some traders also use the term to refer to the sensitivity of an option's theoretical value to a change in volatility (equivalent to the vega)\par
Term Structure The distribution of implied volatilities across different expiration months in the same underlying market.\par
Theoretical Value An option value generated by a mathematical model\par
given certain prior assumptions about the terms of the option, the characteristics of the underlying contract, and prevailing interest rates. Also known as Fair Value.\par

\marksourcepage{824}
Theta (?) The sensitivity of an option's theoretical value to a change in the amount of time remaining to expiration.\par
Three-Way A position similar to a conversion or reversal but where the long or short position in the underlying instrument has been replaced with a very deeply in-the-money call or put.\par
Time Box A long call and short put with the same exercise price and\par
expiration date together with a short call and long put at a different exercise price and expiration date. This is simply a roll using different exercise prices. Also known as a Diagonal Roll.\par
Time Premium See Time Value.\par
Time Value The price of an option less its intrinsic value. The price of an out-of-the-money option consists solely of time value. Also known as Extrinisic Value or Time Premium.\par
Time Spread See Calendar Spread.\par
Type The designation of an option as either a call or a put.\par
Underlying The instrument to be delivered in the event an option  is\par
exercised.\par
Vanilla Option An option, usually exchange traded, with standardized and traditional contract specifications as opposed to an exotic option.\par
Vanna The sensitivity of an option's delta to a change in volatility.\par
Variation The daily cash flow resulting from changes in the settlement price of a futures contract.\par
Vega. The sensitivity of an option's theoretical value to a change in volatility. Also known as Kappa.\par
Vega Decay The sensitivity of an option's vega to the passage of time.\par

\marksourcepage{825}
Vertical Spread The purchase of an option at one exercise price and the sale of an option at a different exercise price where both options are of the same type, have the same underlying contract, and expire at the same time.\par
Volatility The degree to which the price of a contract tends to fluctuate over time.\par
Volatility Skew The tendency of options at different exercise prices to trade at different implied volatilities. Also known as a Volatility Smile.\par
Volatility Smile See Volatility Skew.\par
Volga The sensitivity of an option's vega to a change in volatility. Also known as Vomma.\par
Vomma See Volga.\par
Warrant A long-term call option. The expiration date of a warrant may\par
under some circumstances be extended by the issuer.\par
Write To sell an option.\par
Zero-Cost Collar A collar where the prices of the purchased and sold\par
options are equal.\par
Zomma The sensitivity of an option's gamma to a change in volatility.\par

\chapter{Some Useful Math}

\marksourcepage{826}
\noindent B\par
The mathematical functions and calculations referred to in this text are included in almost all commonly used spreadsheets, and for most traders, it is not necessary to know exactly how the calculations are made. Of far greater importance is the ability to interpret the numbers that result from the calculations.\par
For the reader who is interested, a detailed discussion of these mathematical concepts can be found in any good statistics or finance textbook. For convenience, we include an overview of these concepts and applications.\par
\noindent Rate-of-Return Calculations\par
An interest rate is the most common rate of return. The total interest can be computed in three ways: simple, compound, and continuous. If\par

\marksourcepage{827}
then, for simple interest,\par
for compound interest,\par
and for continuous interest,\par
Because volatility is a continuously compounded rate of return, we can use the exponential and logarithmic functions to do similar calculations for volatility. If\par
\noindent t = time to expiration, in years\par

\marksourcepage{828}
\noindent F = a forward price after the period of time t\par
\noindent sigma = annual volatility or standard deviation\par
\noindent X = an option's exercise price\par
\noindent then a price range of n standard deviations is\par
The number of standard deviations required to reach an exercise price is\par
\noindent Normal Distributions and Standard Deviation\par
\noindent If\par
\noindent xi = each data point\par
\noindent n = number of data points\par
\noindent sigma = standard deviation or volatility\par
\noindent ? = average or mean\par
\noindent then the mean or average ? is\par
When calculating the standard deviation from the entire population, sigma is given by\par
When estimating the standard deviation from a sample of the entire\par
\noindent population, sigma is given by1\par

\marksourcepage{829}
\noindent The normal distribution curve n(x) is given by\par
In a standard normal distribution, ? = 0 and sigma = 1.\par
Many of the measures associated with a distribution are derived from a group of numbers called moments. In general, the jth moment mj about the mean ? of a distribution is given by2\par
From the second, third, and fourth moments, we can calculate the\par
\noindent skewness and kurtosis of a distribution\par
A perfectly normal distribution has a skewness of 0 and a kurtosis of 3. To normalize the kurtosis such that a normal distribution has a kurtosis of 0, it is common to subtract 3\par
Figure B-1 shows calculation of the mean and standard deviation for the pinball distribution in Figure 6-2. The steps required to calculate the skewness and kurtosis would require an inordinate amount of space. However, the relevant values, including the first three moments, are\par

\marksourcepage{830}
(The right tail of the distribution is very slightly longer than the left tail.)\par
(The peak of the distribution is slightly lower and the tails slightly shorter than a true normal distribution.)\par
figure B-1 Calculation of the mean and standard deviation for the distribution in Figure 6-2.\par
\noindent Volatility\par
Volatility is usually calculated as a sample standard deviation. It is also common to assume a mean of 0. The estimated annualized volatility is then\par

\marksourcepage{831}
\noindent given by\par
where xi = ln(pn/pn-1) = natural logarithm of the current price pn divided by the previous price pn-1 and t = the time interval, in years, between price changes.\par
If the underlying contract is a stock, in theory, the price returns xi should be adjusted to reflect the forward price of pn-1 over each time period. However, unless interest rates are very high or the stock will pay a dividend, using the actual price rather than the forward price is unlikely to significantly alter the results.\par
The volatility calculation for the stock option example in Figure 8-1 is shown in Figure B-2. Because price changes were observed at seven-day intervals (t = 7/365), to annualize the volatility, it was necessary to divide by sqrt7/365. The calculation represents the population standard deviation (dividing by n rather than n - 1) and is based on the actual mean of the price changes. We might also calculate the volatility assuming a 0 mean or use an estimated standard deviation. The various results are as follows:\par
There is very little difference between the calculations made from the actual mean and a 0 mean. The estimated standard deviation is always greater than the population standard deviation.\par

\marksourcepage{832}
Figure B-2 Volatility calculation for the stock option example in\par
Figure 8-1.\par

\marksourcepage{833}
1 The sample standard deviation is sometimes denoted with s (instead of sigma). 2 In the same way we calculate a sample standard deviation by dividing by n - 1, we can also calculate sample moments by dividing by n - 1 rather than dividing by n.\par

\chapter{Index}

\marksourcepage{834}
Please note that index links point to page beginnings from the print edition. Locations are approximate in e-readers, and you may need to page down one or more times after clicking a link to get to the indexed material.\par
\noindent Adjusted delta, 499\par
\noindent Adjustments\par
\noindent in Black-Scholes model, 348\par
\noindent dynamic hedging with, 121-122\par
\noindent to original hedge, 123\par
\noindent risks and neutral, 246\par
\noindent to spread, 246-248\par
\noindent trader choices for, 247-248\par
\noindent in volatility spreads, 206-208\par
\noindent All or none (AON), 539\par
\noindent American option, 32, 292, 539\par
\noindent arbitrage boundaries for, 293-295, 300\par
\noindent binomial tree and, 370-371\par
\noindent Black-Scholes model not for, 309, 377\par
\noindent delta values in, 313\par
\noindent dividend-paying stocks and, 373-376\par
\noindent European option compared to, 315-317\par
\noindent evaluating, 314\par
\noindent pricing of, 309-317\par
\noindent AON. See All or none\par
Approximations, ,\par

\marksourcepage{835}
\noindent for put-call-parity, 271-273\par
\noindent for Black-Scholes value, 350-352\par
\noindent Arbitrage, 19-21, 265-266, 539\par
\noindent boundaries, 293-299, 300\par
\noindent free, 59, 60\par
\noindent index, 452-454, 545\par
\noindent relationships, 291\par
\noindent risk, 273-290\par
ARCH. See Autoregressive conditional heteroskedasticity\par
\noindent Asian option. See Average price option\par
\noindent Ask price, 66\par
\noindent Assignment, 29-30, 539\par
\noindent At the forward, 60, 539\par
\noindent At the money, 33-35, 93\par
\noindent bull and bear strategies and, 222-226\par
\noindent delta, 136\par
\noindent gamma, 149-150\par
\noindent straddles, 477\par
\noindent theta, 141\par
\noindent vega, 145-146\par
\noindent Automatic exercise, 35-36, 539\par
Autoregressive conditional heteroskedasticity (ARCH), 394\par
\noindent Average price option, 540\par
\noindent Bachelier, Louis, 61-62, 76\par
\noindent Backspread, 179, 540\par
\noindent Backward, 14, 524-526, 530, 540\par
\noindent Balanced skew, 487\par
\noindent Barone-Adesi, Giovanni, 309\par
\noindent Barone-Adesi-Whaley model, 309-310\par
\noindent Barrier option, 540\par
\noindent Bear spread, 222-226, 540\par

\marksourcepage{836}
\noindent Bermuda option, 540\par
\noindent Bias, in futures market, 457-458\par
\noindent Bid price, 66\par
\noindent Bid-ask spread, 251-252, 427-428, 485\par
\noindent Binary option, 540\par
\noindent Binomial expansion, 361\par
\noindent Binomial model. See Cox-Ross-Rubinstein model\par
\noindent Binomial notation, 363-365\par
\noindent Binomial option pricing, 358\par
\noindent Binomial tree, 359-360\par
\noindent American options and, 370-371\par
\noindent Black-Scholes model values in, 379\par
\noindent call value, 364\par
\noindent dividend-paying stock and, 374-375\par
\noindent with dividends, 372\par
\noindent notation for, 363-365\par
\noindent one-period, 360-361\par
\noindent periods used in, 378\par
\noindent recombining, 368\par
\noindent Black, Fischer, 57, 61, 62\par
\noindent Blackout period, 304\par
\noindent Black-Scholes model, 61-68, 120\par
\noindent approximations of, 350-352\par
\noindent binomial values in, 379\par
\noindent continuous diffusion process assumed in, 472\par
\noindent as continuous-time model, 84\par
\noindent delta in, 352-354\par
\noindent dividends in, 67-68\par
\noindent equation, 339-340\par
\noindent as European pricing model, 309, 377\par
\noindent formulae, 349-350\par
\noindent gamma and vega in, 354-357\par

\marksourcepage{837}
\noindent implied volatility in, 130-131, 307\par
\noindent interest rates in, 66-67\par
\noindent jump-diffusion model variation of, 475\par
lognormal distribution assumptions in, 85, 341-344, 506\par
\noindent not for American options, 309, 377\par
\noindent out of the money options and, 470-471\par
\noindent as probabilistic model, 471\par
\noindent put-call parity in, 340-341\par
\noindent theoretical values calculated with, 63\par
\noindent theta of, 354\par
\noindent time to expiration in, 65-66\par
\noindent underlying price in, 66, 345-346\par
\noindent variations, 62\par
\noindent volatility in, 68, 339\par
\noindent Black-Scholes-Merton model, 338\par
\noindent BMX. See Chicago Board Options Exchange BuyWrite Index\par
\noindent Body of butterfly, 192\par
\noindent Bonds and notes, 17\par
\noindent Borrowing costs rbc, 24\par
\noindent Box, 280-282, 540\par
\noindent Breakeven prices, 51\par
\noindent Breakeven volatility, 116-117, 130\par
\noindent Broad based stock indexes, 441-442\par
\noindent Bull and bear butterfly, 211-212\par
\noindent Bull and bear calendar spread, 212-214\par
\noindent Bull and bear ratio spread, 210-211, 213\par
\noindent Bull and bear vertical spread, 215-217\par
\noindent Bull spread, 222-226, 540\par
\noindent Bull straddle, 171\par
\noindent Bund futures, 388, 390\par
\noindent Butterfly, 289, 506-507, 540\par
\noindent bull and bear, 211-212\par

\marksourcepage{838}
\noindent call or put, 290\par
\noindent iron, 260-262, 290, 545\par
\noindent and implied distributions, 508\par
\noindent long, 174-175, 201\par
\noindent short, 174-175, 201\par
\noindent time, 191-192\par
\noindent Buying, 284, 287\par
\noindent Buy/write, 326-327, 540\par
\noindent Cabinet bid, 420, 540\par
\noindent Calendar spread, 195, 237, 239-240, 540\par
\noindent bull and bear, 212-214\par
\noindent implied volatility and, 405-408\par
\noindent long and short, 189-190, 204, 214\par
\noindent rolls and, 284-286\par
\noindent Call\par
\noindent butterfly, 290\par
\noindent Christmas tree, 183, 203\par
\noindent covered, 486\par
\noindent delta, 137, 139\par
\noindent lambda and, 154-156\par
\noindent option, 3, 26, 299-302, 540-541\par
\noindent protective, 322-324\par
\noindent ratio spread, 179-181, 203\par
\noindent synthetics, 288, 551\par
\noindent theoretical value of, 101, 108-109, 111, 310-311\par
\noindent Cap, 322, 514, 540, 541\par
\noindent Capitalization, 443\par
\noindent Capitalization weighted index, 443, 454-456\par
\noindent Carry trade, 20\par
\noindent Cash, 1, 8, 31-32, 459-461, 461-462\par
\noindent Cash-and-carry arbitrage, 20\par

\marksourcepage{839}
\noindent Cash-and-carry strategy, 159, 161\par
\noindent Cash-secured put, 327\par
\noindent Castelli, Charles, 61\par
\noindent CBOE. See Chicago Board Options Exchange\par
\noindent CEV. See Constant elasticity of variance model\par
\noindent Charm, 139-141, 541\par
Chicago Board Options Exchange (CBOE), 62, 250, 309, 459, 515, 530\par
Chicago Board Options Exchange BuyWrite Index (BXM), 327\par
\noindent Chicago Mercantile Exchange, 36, 67, 81, 452\par
\noindent Chooser option, 541\par
\noindent Christmas tree, 182-184, 203, 541\par
\noindent Circuit breakers, 464\par
\noindent Clearing firm, 11\par
\noindent Clearing member, 541\par
\noindent Clearinghouse, 10, 11, 541\par
\noindent Closing trade, 5\par
\noindent CME Clearing House, 10\par
\noindent Collar, 328-330, 541\par
\noindent Color, of option, 152-153, 541\par
\noindent Combination (Combo), 266, 541\par
\noindent Combo value, 266, 267\par
\noindent Commodity markets, 486\par
\noindent Companion option, 256, 259\par
\noindent Complex positions, 45-46, 410-411, 420-423\par
\noindent Compound option, 541\par
\noindent Condor, 176-178, 541\par
\noindent Constant elasticity of variance (CEV) model, 478\par
\noindent Contango, 14, 524-526, 530, 541\par
\noindent Contingency order, 208, 542\par
\noindent Continuous diffusion process, 471-472\par
\noindent Continuous-time model, 84\par
\noindent Contract specifications, 26-30\par

\marksourcepage{840}
\noindent Convenience yield, 15\par
\noindent Conversion market, 267\par
\noindent Conversion, 265-266, 273, 275, 277-282, 289, 542\par
\noindent Counterparty, 10\par
\noindent Covered call, 486\par
\noindent Covered option, 329\par
\noindent Covered write, 61, 324-328, 542\par
\noindent Cox, John, 309, 358\par
\noindent Cox-Ross-Rubinstein model, 309-310\par
\noindent Crack spread, 163-164\par
\noindent Credit spread, 214\par
\noindent Curvature, 105-108\par
\noindent Cylinder. See Collars\par
\noindent DAX Index, 448\par
\noindent Debit spread, 214\par
\noindent Declared date, 22-23\par
\noindent Delta (Delta), 64, 110, 492, 542\par
\noindent adjusted, 499\par
\noindent of American options, 313\par
\noindent in Black-Scholes model, 352-354\par
\noindent of bull spread, 218-219, 223\par
\noindent call, 137, 139\par
\noindent decay, 139\par
\noindent hedge ratio and, 102-103\par
\noindent implied, 136\par
\noindent at the money, 136\par
\noindent negative and positive, 416\par
\noindent position, 169\par
\noindent probability and, 104-106\par
\noindent put, 138\par
\noindent rate of change and, 100-102\par

\marksourcepage{841}
\noindent theoretical underlying position and, 103-104\par
\noindent as volatility changes, 135-136\par
\noindent Delta (directional) risk, 227\par
Delta-neutral, 103, 121, 129-130, 133, 353-354, 421, 542\par
\noindent Derivative contracts, 4\par
\noindent Dermna, Emanuel, 57\par
\noindent Diagonal ratio spread, 236-237\par
\noindent Diagonal roll. See Time box\par
\noindent Diagonal spread, 196-200, 542\par
\noindent Diffusion process, 472-473\par
\noindent Digital option. See Binary option\par
\noindent Dividend play, 318\par
\noindent Dividend value, 300, 302\par
\noindent Dividend-paying stocks, 373-376\par
\noindent Dividends, 279, 373-379\par
\noindent binomial tree with, 372\par
\noindent in Black-Scholes model, 67-68\par
\noindent declared date of, 22-23\par
\noindent Dow Jones Industrial Average paying, 452\par
\noindent ex-date of, 22-23, 68, 543\par
\noindent implied, 21, 273\par
\noindent payable date of, 23\par
\noindent record date of, 22\par
\noindent risk and, 239-240, 278-280\par
\noindent stock options influenced by, 100\par
\noindent calendar spreads and, 192-195\par
\noindent Divisor, 446-447\par
\noindent Dow Jones Industrial Average, 445, 452\par
\noindent Downside contract position, 418\par
\noindent Dragonfly, 504, 542\par
\noindent Driftless theta, 354\par
\noindent Dynamic hedging, 336, 474, 512, 542\par

\marksourcepage{842}
\noindent with adjustments, 121-122\par
\noindent delta-neutral, 133\par
\noindent process of, 127-129\par
\noindent Early exercise, 292\par
\noindent of call option, 299-302\par
\noindent of future option, 305-307\par
\noindent protective value and, 308\par
\noindent of put option, 302-305\par
\noindent risk, 319-320\par
\noindent short stock impact on, 305\par
\noindent strategies, 317-319\par
\noindent Efficiency, 244-246, 542\par
\noindent Efficient-market hypothesis, 394\par
\noindent Elasticity, 154, 542. See also Lambda\par
\noindent Equal-weighted index, 444\par
\noindent Equilibrium price, 429\par
\noindent Equity portfolio, 457\par
\noindent Euro-bund options, 269\par
\noindent Eurocurrency, 81-82, 542\par
\noindent Eurocurrency rate, 543\par
\noindent European box, 314\par
\noindent European option, 32, 60, 62, 253, 292, 543\par
\noindent American option compared to, 315-317\par
\noindent arbitrage boundaries for, 293-299\par
\noindent arbitrage relationships for, 291\par
\noindent Black-Scholes model pricing of, 309, 377\par
\noindent dividend-paying stocks and, 373\par
\noindent European pricing model, 309, 377\par
\noindent EuroStoxx 50 Index, 385-386, 402, 405, 407\par
\noindent EWMA. See Exponentially weighted moving average\par
\noindent Exchange option, 543\par

\marksourcepage{843}
\noindent Exchanges, 6-9, 449, 464\par
\noindent Exchange-traded contracts, 10, 160\par
\noindent Exchange-traded funds, 454\par
\noindent Exchange-traded options, 29\par
\noindent Ex-dividend date (Ex-date), 22-23, 68, 543\par
\noindent Execution risk, 273-274\par
\noindent Exercise, 29-30, 543\par
\noindent Exercise notice, 35\par
\noindent Exercise price, 4, 29, 65, 215, 543\par
\noindent Exotic option, 182, 543\par
\noindent Expected value, 53-54, 59-60\par
\noindent Expiration date, 2, 4, 28-29, 65-66, 73, 543\par
\noindent Expiration straddle, 476-477\par
\noindent Expiry. See Expiration date\par
\noindent Exponentially weighted moving average (EWMA), 393-394\par
\noindent Extreme-value method, 384\par
\noindent Extrinsic value. See Time value\par
\noindent Fair value. See Theoretical value\par
\noindent Fear index, 523\par
\noindent Fence. See Collars\par
\noindent Figlewski, Stephen, 57\par
\noindent Fill or kill (FOK), 208, 543\par
\noindent Fixed-income markets, 161-162\par
\noindent Flat peak (platykurtic), 481\par
\noindent Flat skew, 487\par
\noindent Flex option, 27, 543\par
\noindent Floating skew, 489, 490\par
\noindent Floor, 322-323, 543\par
\noindent FOK. See Fill or kill\par
\noindent Forecasting, volatility, 391-394\par
\noindent Foreign currencies, 17-18\par

\marksourcepage{844}
\noindent Foreign currency options, 99\par
\noindent Forward contract, 2, 543\par
\noindent basis, 13\par
\noindent fair price for, 12-14\par
\noindent for foreign currency, 17-18\par
\noindent interest rate risk, 21\par
\noindent long-term, 452\par
\noindent notional value, 6\par
\noindent stock index, 455\par
\noindent synthetics, 254\par
\noindent Forward conversion. See Conversion\par
\noindent Forward price, 2-3, 12-17, 76-77, 270-271, 543\par
\noindent Forward rate, 16, 408\par
\noindent Forward start option, 544\par
\noindent Forward volatility, 404-409\par
\noindent Forward-rate agreement (FRA), 16\par
\noindent Free float, 445\par
\noindent Frequency distribution, 480-481\par
\noindent Frictionless markets, 125-126, 463-466\par
\noindent Front spread, 182, 544\par
\noindent FTSE 100 index, 485\par
\noindent Fugit, 307, 544\par
\noindent Future volatility, 394-397\par
\noindent Futures contract, 2, 544\par
\noindent early exercise of, 305-307\par
\noindent locked markets, 269-270\par
\noindent notional value, 455\par
\noindent option's difference with, 26-27\par
\noindent options on, 267-269\par
\noindent on stock index, 450-451\par
\noindent synthetics, 270, 277\par
\noindent Futures exchange, 2\par

\marksourcepage{845}
\noindent Futures market, 457-458, 464\par
\noindent Futures option, 19, 30-31\par
\noindent Futures-type settlement, 7-9, 36, 112, 544\par
\noindent Gamma (?), 110, 113, 152, 544\par
\noindent in Black-Scholes model, 354-357\par
\noindent of bull spreads, 223\par
\noindent influence of expiration and volatility, 150-151\par
\noindent magnitude risk measured by, 115-116\par
\noindent at the money, 149-150\par
\noindent negative and positive, 106-107, 230, 414-418, 421\par
\noindent options, 154, 367\par
\noindent as option's curvature, 105-108\par
\noindent rent, 369-370\par
\noindent spread, 165\par
\noindent theta tradeoff with, 237\par
\noindent underlying price changes and, 149-150\par
\noindent Gamma (curvature) risk, 227-228\par
\noindent Gamma spread, 165\par
\noindent Gaps, real world, 476\par
GARCH. See Generalized Autoregressive conditional heteroskedasticity\par
\noindent Garman, Mark, 62, 385\par
\noindent Garman-Klass estimators, 385\par
\noindent Garman-Kohlhagen model, 62, 67\par
Generalized Autoregressive conditional heteroskedasticity (GARCH), 394 Geometric-weighted index, 445 Go long or short, 5 Gold, 93-96, 384, 387, 389 Good `til canceled (GTC), 544 Greeks, 99 GTC. See Good `til canceled Guts, 172, 544\par

\marksourcepage{846}
\noindent Haircut, 8, 544\par
\noindent Half-steps, 378\par
Hedge ratio. See Delta (?)\par
\noindent Hedgers, 321, 544\par
\noindent Hedging\par
\noindent collars used in, 328-330\par
\noindent complex strategies in, 331-333\par
\noindent covered-writes used in, 324-328\par
\noindent high and low implied volatility in, 331\par
\noindent natural long or short in, 321\par
\noindent protective calls and puts in, 322-324\par
\noindent stock index futures and, 457\par
\noindent strategy summary of, 330\par
\noindent using vertical spreads, 332-333\par
\noindent using volatility spreads, 332\par
\noindent volatility contracts used for, 536\par
\noindent volatility reduced by, 333-334\par
\noindent Higher-order risk measures, 355\par
\noindent Historical volatility, 381-386, 392\par
\noindent of Bund futures, 388, 390\par
\noindent of EuroStoxx 50, 386, 402\par
\noindent of gold, 93-96, 384, 389\par
\noindent of S\&P 500 Index, 87, 383, 389\par
\noindent Hockey-stick diagrams, 40\par
\noindent Horizontal spread. See Calendar spread\par
\noindent Illiquid market, 66\par
\noindent Immediate or cancel (IOC), 544\par
\noindent Implied delta, 136\par
\noindent Implied distribution, 506-510\par
\noindent Implied dividend, 21, 273\par
\noindent Implied interest rates, 21, 273\par

\marksourcepage{847}
\noindent Implied spot price, 21\par
\noindent Implied volatility, 88-89, 92-95, 116, 199, 544\par
\noindent assessing risk of, 401\par
\noindent in Black-Scholes model, 130-131, 307\par
\noindent calendar spread and, 405-408\par
\noindent contracts, 512, 514-523\par
\noindent defining, 514\par
\noindent EuroStoxx 50 and, 407\par
\noindent floating skews as, 490\par
\noindent future volatility predicted by, 394-397\par
\noindent hedging with high and low, 331\par
\noindent influence of, 188-189\par
\noindent option pricing compared with, 90-91\par
\noindent reevaluation of, 131\par
\noindent relative changes in, 400\par
\noindent sensitivity to changes in, 425-426\par
\noindent term structure of, 397-404, 430-432\par
\noindent time decay and, 380-381\par
\noindent In price, 544\par
\noindent In the money, 33-35, 92, 545\par
\noindent Index\par
\noindent arbitrage, 452-454, 545\par
\noindent broad based stock, 441-442\par
\noindent calculation, 443-445\par
\noindent capitalization-weighted, 443, 454-456\par
\noindent divisor, 446-447\par
\noindent equal-weighted, 444\par
\noindent fear, 523\par
\noindent geometric-weighted, 445\par
\noindent impact of individual stock price changes, 448-450\par
\noindent narrow based stock, 441-442\par
\noindent price-weighted, 443-445, 449, 454-455\par

\marksourcepage{848}
\noindent rebalancing, 445\par
\noindent replicating, 454-456\par
\noindent total-return, 447-448\par
\noindent volume-weighted average price, 450\par
\noindent In-option, 544\par
\noindent In-price, 544\par
\noindent Insurance policies, 4, 321\par
\noindent Interest, 5, 123-124, 239-240, 278-280, 303, 318-319\par
\noindent Interest rates, 228\par
\noindent in Black-Scholes model, 66-67\par
\noindent foreign currency options and, 99\par
\noindent forward contract and risks of, 21\par
\noindent implied, 21, 273\par
\noindent negative, 83\par
\noindent option values influenced by, 97-100, 466-467\par
\noindent products volatility of, 81-82\par
\noindent stock options importance of, 305-306, 466-467\par
\noindent in theoretical pricing models, 126\par
\noindent theoretical value and, 467\par
\noindent volatility spreads with dividends and, 192-195\par
\noindent Interest-rate value, 300\par
\noindent Intermarket spreads, 161-162, 545\par
\noindent Intrinsic value, 32-34, 38-40, 299-300, 545\par
\noindent Investment skew, 486-487\par
\noindent IOC. See Immediate or cancel\par
\noindent Iron butterfly, 260-262, 290, 545\par
\noindent Iron condor, 262-263, 545\par
\noindent Jelly roll. See Roll\par
\noindent Jump process, 472, 473\par
\noindent Jump-diffusion model, 475, 478\par
\noindent Jump-diffusion process, 472, 473\par

\marksourcepage{849}
\noindent Kansas City Board of Trade, 452\par
\noindent Kappa (K), 110, 545\par
\noindent Klass, Michael, 385\par
\noindent Knock-in option. See In-option\par
\noindent Knock-out option. See Out-option\par
\noindent Kohlhagen, Steven, 62\par
\noindent Kurtosis, 481-483, 495-498, 501-506, 555\par
\noindent Ladder, 545. See also Christmas tree\par
\noindent Lambda (Lambda), 154-156. See also Elasticity\par
\noindent Last trading day, 28\par
\noindent LEAP. See Long-term Equity Anticipation Security\par
\noindent Leg, 163-164, 545\par
\noindent Leverage value. See Elasticity\par
\noindent LIBOR. See\par
\noindent London Interbank Offered Rate\par
\noindent Limit, 546\par
\noindent Limit order, 208, 546\par
\noindent Liquid market, 66\par
\noindent Liquidity, 249-250, 427\par
\noindent Local, 546\par
\noindent Locked limit, 23\par
\noindent Locked markets, 23, 546\par
\noindent Lognormal distributions, 82-83\par
Black-Scholes model and assumptions of, 85, 341-344, 506\par
\noindent butterflies with, 508\par
\noindent price changes in, 84\par
\noindent standard deviations and, 346-348\par
\noindent of underlying price, 478-481\par
\noindent London Interbank Offered Rate (LIBOR), 67, 81\par
\noindent Long, 54, 321, 477\par
\noindent butterfly, 174-175, 201\par

\marksourcepage{850}
\noindent calendar spread, 189-190, 204, 214\par
\noindent call, 38, 42-43\par
\noindent call Christmas tree, 183, 204\par
\noindent condor, 177-178, 202\par
\noindent natural, 321\par
\noindent position, 5, 98, 328, 546\par
\noindent premium, 197, 546\par
\noindent put, 39, 42-43\par
\noindent put Christmas tree, 184, 203\par
\noindent rate rl, 24\par
\noindent ratio spread, 546\par
\noindent straddle, 170-171, 205\par
\noindent strangle, 172-173, 201, 262\par
\noindent synthetic, 253, 260-261\par
\noindent time butterfly, 193\par
\noindent Long-term Equity Anticipation Security (LEAP), 545\par
\noindent Long-term options, 27-28\par
\noindent Lookback option, 546\par
\noindent Magnitude risk, 115-116\par
\noindent Margin and variation settlement, 7\par
\noindent Margin deposit, 7\par
\noindent Margin of error, 239, 240-244\par
\noindent Margins, 8, 546\par
\noindent Market conditions\par
\noindent market making and, 429-430\par
\noindent option traders and changes in, 139-140\par
\noindent option values influenced by, 97-98\par
\noindent positions influenced by, 435-436\par
\noindent risk analysis and changes in, 434-435, 437\par
\noindent risk and, 413, 415-416\par
\noindent stock splits and, 438-440\par

\marksourcepage{851}
\noindent theoretical edge and, 231\par
\noindent Market integrity, 10-11\par
\noindent Market liquidity, 235\par
\noindent Market maker, 427-430, 432-433, 546\par
\noindent Market order, 208, 546\par
\noindent Market-if-touched (MIT), 546\par
\noindent Market-on-close (MOC), 453, 546\par
\noindent Markets, 126, 165-166\par
\noindent commodity, 486\par
\noindent declining, 426-427\par
\noindent fixed-income, 161-162\par
\noindent frictionless, 463-466\par
\noindent futures contracts and locked, 23, 269-270, 546\par
\noindent option position in, 64\par
\noindent option traders and speed of, 53\par
\noindent volatility measuring speed of, 69\par
\noindent Mark-to-market, 546\par
\noindent Married put, 547\par
\noindent Maturity date. See expiration date\par
\noindent Mean, 73-77, 345\par
\noindent Mean reverting, 387-388, 398\par
\noindent Median, 345\par
\noindent Merton, Robert, 62, 338\par
\noindent Mid-Atlantic option. See Bermuda option\par
\noindent Midcurve options, 28, 547\par
\noindent MIT. See Market-if-touched\par
\noindent MOC. See Market-on-close\par
\noindent Models. See specific model\par
\noindent Modes, 84, 345\par
\noindent Moneyness, 491\par
\noindent Naked position, 164, 209, 547\par

\marksourcepage{852}
\noindent Narrow based stock index, 441-442\par
\noindent Natural gas futures, 403-404\par
\noindent Natural longs and shorts, 321\par
\noindent Negative dividend risk, 279\par
\noindent Negative gamma position, 106-107, 230, 414-418, 421\par
\noindent Negative interest rates, 83\par
\noindent Negative theoretical edge, 118\par
\noindent Negative time value, 109-110\par
\noindent Negative vega, 230, 236\par
\noindent Negative volga, 236\par
\noindent Net contract position, 417\par
\noindent Neutral hedge. See Riskless hedge\par
\noindent Neutral spread, 210, 547\par
\noindent New York Mercantile Exchange, 29\par
\noindent Nominal value. See notional value\par
\noindent Nonsymmetrical strategies, 185\par
\noindent Normal distribution, 69-73, 75-77, 480-481, 554-556\par
\noindent Not held, 209, 547\par
\noindent Notional value, 6, 455\par
\noindent Notional vega, 513\par
\noindent OBO. See Order book official\par
\noindent OCO. See One-cancels-the-other\par
\noindent OEX. See Options Exchange Index Omega (?), 154, 547\par
\noindent One-cancels-the-other (OCO), 209, 547\par
\noindent Open interest, 5\par
\noindent Open position, 5\par
\noindent Opening trade, 5\par
\noindent Option pricing, 32-36, 62, 90-91, 338-339, 358\par
\noindent Option replication. See Portfolio insurance\par
\noindent Option risk analysis, 116\par
\noindent Option-pricing theory, 119\par

\marksourcepage{853}
\noindent Options Clearing Corporation, 10, 36\par
\noindent Options Exchange Index (OEX), 459, 515\par
\noindent Order book official (OBO), 547\par
\noindent Orders, submitting spread, 208-209\par
\noindent OTC. See Over-the-counter market\par
Out of the money, 33-35, 93, 135-136, 419-420, 423, 470-471, 547\par
\noindent Outliers, 481\par
\noindent Out-option, 547\par
\noindent Out-price, 547\par
\noindent Out-trade, 547\par
\noindent Over-the-counter (OTC) market, 19\par
\noindent Overwrite, 324, 547\par
\noindent Paralysis through analysis, 118\par
\noindent Parity. See Intrinsic value\par
\noindent Parity graphs, 38-51\par
\noindent Parkinson, Michael, 384\par
\noindent Partial derivatives, 100\par
\noindent Path dependent, 471\par
\noindent Phi (Phi), 112\par
\noindent Physical commodities, 2, 14-15\par
\noindent Physical settlement, 7-8\par
\noindent Physical underlying, 30-31\par
\noindent Pin risk, 274-275\par
\noindent P\&L. See Profit and loss\par
\noindent Plain-vanilla-interest-rate swap, 5\par
\noindent Population standard deviation, 382\par
\noindent Portfolio insurance, 335-337\par
\noindent Portfolio manager, 333-334, 336, 537\par
\noindent Positive dividend risk, 279\par
\noindent Positive gamma position, 106-107, 415\par
\noindent Premium, 4\par

\marksourcepage{854}
\noindent Price changes, 479\par
\noindent individual stock, 448-450\par
\noindent in lognormal distributions, 84\par
\noindent underlying, 141-142, 145-146, 149-150\par
\noindent volatility and observed, 80-81\par
\noindent Price distribution, 73\par
\noindent Price movement, 65, 72-75\par
\noindent Price volatility, 82\par
\noindent Price weighted index, 443-445, 449, 454-455\par
\noindent Pricing, of American options, 309-317\par
\noindent Pricing models, 62, 338-339\par
\noindent American, 309\par
\noindent entering skew into, 489-494\par
\noindent European, 309, 377\par
\noindent volatility skew in, 499\par
\noindent weaknesses, 484-485\par
\noindent Probabilistic model, 471\par
\noindent Probabilities, 4, 376, 483\par
\noindent delta and, 104-106\par
\noindent expected value, 53-54, 59-60\par
\noindent models and, 56-57\par
\noindent symmetrical distribution of, 60\par
\noindent theoretical values in, 54-56\par
\noindent Probability functions, 344\par
\noindent Profit and loss (P\&L), 46-51, 124, 125, 133\par
\noindent Program trading, 453\par
\noindent Protective options, 323-324, 327, 329\par
\noindent Protective put, 486\par
\noindent Protective value, 308\par
\noindent Pseudoprobabilities, 376\par
\noindent Put\par
\noindent butterfly, 290\par

\marksourcepage{855}
\noindent cash-secured, 327\par
\noindent Christmas tree, 184, 203\par
\noindent delta, 138\par
\noindent lambda and, 154-156\par
\noindent long, 39, 42-43\par
\noindent married, 547\par
\noindent option, 4, 26, 302-305\par
\noindent protective, 322-324, 486\par
\noindent ratio spread, 179-181, 204\par
\noindent short, 40\par
\noindent short Christmas tree, 184\par
\noindent synthetics, 261-262, 551\par
\noindent theoretical value of, 101-102, 108-109, 111, 312\par
\noindent Put-call parity, 267-268, 462\par
\noindent in Black-Scholes model, 340-341\par
\noindent for stock options, 270-273, 288\par
\noindent synthetic futures contract and, 270\par
\noindent Quadratic model. See Barone-AdesiWhaley model\par
\noindent Random walk, 69-73, 471\par
\noindent Rate of return, 82, 553-554\par
Ratio spread, 179-182, 202, 210-211, 213, 236-237, 546, 549\par
\noindent Ratio straddle, 171\par
\noindent Ratio strategy, 162\par
\noindent Ratio write, 332\par
\noindent Realized profit, 7\par
\noindent Realized volatility, 86-87, 116, 512-514\par
\noindent Rebalancing process, 537\par
\noindent Recombining binomial tree, 368\par
\noindent Record date, 22\par
\noindent Rehedging, 129, 367\par

\marksourcepage{856}
\noindent Reversal, 273, 275, 277-282, 502, 549\par
\noindent Reverse conversion, 265, 268, 278-279, 289\par
\noindent Rho (P), 111-112, 113, 467, 549\par
\noindent Rho (interest-rate) risk, 228\par
\noindent Risk, 21, 273-274\par
\noindent analysis, 434-435, 437\par
\noindent arbitrage, 273-290\par
\noindent boxes, 280-282\par
\noindent complex positions, 422-423\par
\noindent delta (directional), 227\par
\noindent dividends and interest, 239-240, 278-280\par
\noindent early exercise, 319-320\par
\noindent gamma (curvature), 227-228\par
\noindent gamma spreads and, 236-237\par
\noindent interest and, 239-240, 278-280\par
\noindent magnitude, 115-116\par
\noindent management, 97-100, 135\par
\noindent margin for error in, 239\par
\noindent negative and positive dividend, 279\par
\noindent neutral adjustments, 246\par
\noindent pin, 274-275\par
\noindent rho (interest-rate), 228\par
\noindent rolls and, 282-285\par
\noindent settlement, 276-278\par
\noindent spreading strategies and, 161, 166-168\par
\noindent in stock options, 240\par
\noindent strategies and considerations of, 234, 238\par
\noindent theoretical pricing models and types of, 463\par
\noindent theta (time decay), 228\par
\noindent time boxes and, 285-287\par
\noindent vega (volatility), 228\par
\noindent volatility, 228-234\par

\marksourcepage{857}
\noindent Risk measures, 99\par
\noindent higher-order, 355\par
\noindent interpreting, 112-118\par
\noindent skewed, 498-499\par
\noindent stock splits and, 439-440\par
\noindent traditional and nontraditional, 156-157\par
\noindent Risk reversal, 502, 549\par
\noindent Risk-analysis process, 432\par
\noindent Risk-free rate, 67\par
\noindent Riskless hedge, 63-64, 102\par
\noindent Risk-neutral world, 358-360\par
\noindent Risk-reward tradeoff, 40, 97, 176, 244\par
\noindent Roll, 282-286, 549\par
\noindent Ross, Stephen, 358\par
\noindent Rubinstein, Mark, 358\par
\noindent Sample standard deviation, 382\par
\noindent Scalper, 158, 549\par
\noindent Scholes, Myron, 61, 62, 338\par
\noindent Second-generation option. See Exotic option\par
\noindent Sell stock short, 23\par
\noindent Serial correlated, 386\par
\noindent Serial options, 27, 549\par
\noindent Settlement in to cash, 31-32\par
\noindent dates, 22\par
\noindent in to a futures position, 30-31\par
\noindent futures-type, 7-9, 36, 112, 544\par
\noindent margin and variation, 7\par
\noindent options and types of, 36\par
\noindent physical, 7-8\par
\noindent in to the physical underlying, 30-31\par
\noindent procedures, 6-9\par

\marksourcepage{858}
\noindent risk, 276-278\par
\noindent stock-type, 7-9, 36, 268, 316, 318-319, 550\par
\noindent Sharpe, William, 334\par
\noindent Sharpe ratio, 334\par
\noindent Short, 5\par
\noindent butterfly, 174-175, 201\par
\noindent calendar spread, 189-190, 204, 214\par
\noindent call, 39\par
\noindent call Christmas tree, 183, 203\par
\noindent collars, 330\par
\noindent condor, 177-178, 202\par
\noindent natural long or, 321\par
\noindent position, 5, 98, 328, 549\par
\noindent premium, 197, 549\par
\noindent put, 40\par
\noindent put Christmas tree, 184, 203\par
\noindent sales, 23-25\par
\noindent sell stock, 23\par
\noindent squeeze, 23, 195, 549\par
\noindent stock, 98, 305\par
\noindent straddle, 170-171, 201\par
\noindent strangle, 172-173, 201, 262\par
\noindent Short rate rs, 24\par
\noindent Short ratio spread, 549\par
\noindent Short-stock rebate, 24\par
\noindent Short-stock squeeze, 23, 195, 549\par
\noindent Short-term options, 27-28\par
\noindent Sigma (sigma), 77, 549\par
\noindent Skew\par
\noindent balanced, 487\par
\noindent flat, 487\par
\noindent floating, 489, 490\par

\marksourcepage{859}
\noindent investment, 486-487\par
\noindent model input of, 498\par
\noindent modeling, 489-494\par
\noindent risk measures and, 498-499\par
\noindent sticky-delta, 492\par
\noindent sticky-strike, 489\par
\noindent volatility, 485-486, 532, 552\par
\noindent Skewness, 481-482, 495-498, 501-506, 555\par
\noindent Slope, in parity graphs, 40-46\par
\noindent S\&P 500 Index, 250, 327, 448, 478-479, 482, 515\par
\noindent historical volatility, 87, 383, 389\par
\noindent SPAN. See Standard Portfolio Analysis of Risk\par
\noindent Special opening rotation, 520\par
\noindent Specialist, 549-550\par
\noindent Speculator, 53, 131, 321, 410, 428, 550\par
\noindent Speed, 150-151, 550\par
Split-strike conversion, 328, 550. See also Risk reversal\par
\noindent Spread, 158, 550\par
\noindent adjustments to, 246-248\par
\noindent bear, 222-226, 540\par
\noindent bid-ask, 251-252\par
\noindent bull, 222-226, 540\par
\noindent common features of, 169\par
\noindent credit and debit, 214\par
\noindent diagonal, 196-200, 542\par
\noindent diagonal ratio, 236-237\par
\noindent dynamic hedging and, 165\par
\noindent efficiency of, 244-246\par
\noindent front, 182, 544\par
\noindent for futures contracts, 160-161\par
\noindent gamma, 165\par
\noindent gamma risk of, 236-237\par

\marksourcepage{860}
\noindent intermarket and intramarket, 161-162\par
\noindent liquidity in, 249-250\par
\noindent with negative vega, 230\par
\noindent static, 165\par
\noindent submitting order for, 208-209\par
\noindent theoretical value and deltas in, 218-219\par
\noindent trader buying, 164\par
\noindent trading style and, 248-249\par
\noindent volatility characteristics and, 233-234\par
\noindent Spreading strategies\par
\noindent defining, 159-164\par
\noindent futures contracts using, 160-161\par
\noindent in options markets, 165-166\par
\noindent risk controlled by, 166-168\par
\noindent risks reduced through, 161\par
\noindent synthetics used in, 259-260\par
\noindent Standard cumulative normal distribution function, 344\par
\noindent Standard deviation\par
\noindent Black-Scholes model, 344-346\par
\noindent calculating, 382\par
\noindent lognormal distributions and, 346-348\par
\noindent mean and, 73-76\par
\noindent normal distribution and, 554-556\par
\noindent price movement and, 74-75\par
\noindent sample and population, 382\par
\noindent volatility as, 77-78\par
\noindent weekly, 79\par
\noindent Standard Portfolio Analysis of Risk (SPAN), 36\par
\noindent Static hedge, 121\par
\noindent Static spread, 165\par
\noindent Sticky-delta skew, 492\par
\noindent Sticky-strike skew, 489\par

\marksourcepage{861}
\noindent Stock index\par
\noindent broad and narrow based, 441-442\par
\noindent forward contract, 455\par
\noindent futures, 458-459\par
\noindent futures contract on, 450-451\par
\noindent hedging and futures in, 457\par
\noindent implied distribution of, 509\par
\noindent options, 458-462\par
\noindent Stock options, 19\par
\noindent dividends influencing, 100\par
\noindent interest rates important in, 305-306, 466-467\par
\noindent long and short positions, 98\par
\noindent physical underlying and, 30-31\par
\noindent put-call parity for, 270-273, 288\par
\noindent Stock splits, 438-440\par
\noindent Stocks\par
\noindent dividend-paying, 373-376\par
\noindent forward price for, 15-17\par
\noindent Index impacted by price changes in, 448-450\par
\noindent Stock-type settlement, 7-9, 36, 268, 316, 318-319, 550\par
\noindent Stop (loss) order, 209, 550\par
\noindent Stop-limit order, 209, 550\par
\noindent Straddles, 61, 170-171, 205, 238, 474, 550\par
\noindent buying, 287\par
\noindent delta neutral and, 353-354\par
\noindent expiration, 476-477\par
\noindent at the money, 477\par
\noindent synthetics and, 260-261\par
\noindent Strangles, 171-173, 238, 261, 550\par
\noindent Strap, 550\par
\noindent Strategies\par
\noindent bull and bear, 222-226\par

\marksourcepage{862}
\noindent buy/write, 326-327, 540\par
\noindent cash-and-carry, 159, 161\par
\noindent choosing appropriate, 199-206\par
\noindent complex hedging, 331-333\par
\noindent early exercise, 317-319\par
\noindent efficiency of, 244-246\par
\noindent hedging summary of, 330\par
\noindent nonsymmetrical, 185\par
\noindent ratio, 162\par
\noindent risk considerations in, 234, 238\par
\noindent skewness and kurtosis, 501-506\par
\noindent symmetrical, 178\par
\noindent Strike price. See Exercise price\par
\noindent Strip, 519, 534-536, 550\par
\noindent Swap, 5, 550\par
\noindent Swaption, 550\par
\noindent Symmetrical strategies, 178\par
\noindent Synthetics, 287-290, 550\par
\noindent call, 288, 551\par
\noindent contract types, 258\par
\noindent forward contract, 254\par
\noindent futures contract, 270, 277\par
\noindent long straddles, 260-261\par
\noindent long underlying, 253\par
\noindent options, 255-260\par
\noindent put, 261-262, 551\par
\noindent relationships, 267\par
\noindent short position, 264-265\par
\noindent spreading strategies using, 259-260\par
\noindent underlying, 253-255, 551\par
\noindent volatility spread using, 287-290\par

\marksourcepage{863}
\noindent Tailing, 277\par
\noindent Tall peak (leptokurtic), 481\par
\noindent Tau (tau), 551\par
\noindent Taxes, 127, 465\par
\noindent Term structure, 390, 397-405, 430-432, 551\par
\noindent Theoretical edge, 118, 227\par
\noindent margin of error and, 240-244\par
\noindent market conditions and, 231\par
\noindent market liquidity and, 235\par
\noindent Theoretical pricing models basic inputs in, 335\par
\noindent diffusion process assumed in, 472-473\par
\noindent interest rates in, 126\par
\noindent performance of, 119-121\par
\noindent problems with, 56-57\par
\noindent risk types in, 463\par
\noindent underlying price in, 66\par
\noindent understanding and having faith in, 538\par
\noindent volatility in, 80, 89-90, 204, 206, 380\par
\noindent weaknesses in, 511-512\par
\noindent Theoretical value, 142-143, 551\par
\noindent Black-Scholes model calculating, 63\par
\noindent of call, 101, 108-109, 111, 310-311\par
\noindent in dynamic hedging, 133-134\par
\noindent of futures contract, 316\par
\noindent interest rates and, 467\par
\noindent of options, 55, 58\par
\noindent and probability, 54-56\par
\noindent of put, 101-102, 108-109, 111, 312\par
\noindent of underlying contract, 115\par
\noindent volatility changes and, 146-147\par
\noindent ``The Theory of Options in Stocks and Shares,'' 61\par
\noindent The Theory of Speculation (Bachelier), 61-62\par

\marksourcepage{864}
\noindent Theta (?), 113, 142-143, 551\par
\noindent of Black-Scholes model, 354\par
\noindent of bull spreads, 224-225\par
\noindent driftless, 354\par
\noindent gamma tradeoff with, 237\par
\noindent at the money, 141\par
\noindent negative time value and, 109-110\par
\noindent of options, 367-368\par
\noindent positive, 417\par
\noindent as time decay, 108-109\par
\noindent underlying price changes and, 141-142\par
\noindent volatility changes and, 144-145\par
\noindent Theta (time decay) risk, 228\par
\noindent Three-way, 280, 551\par
\noindent Ticker symbols, 460\par
\noindent Time, volatility scaled for, 78-79\par
\noindent Time box, 285-287, 551\par
\noindent Time butterfly, 191-192\par
\noindent Time decay, 108-109, 380-381\par
\noindent Time premium, 33, 551\par
\noindent Time spreads. See Calendar spread\par
\noindent Time value, 32, 33-34, 543, 551\par
\noindent Time-series analysis, 393\par
\noindent Total-return index, 447-448\par
\noindent Trading, exchanges halting, 449, 464\par
\noindent Transaction costs, 126-127, 129\par
\noindent Type, 6, 551\par
\noindent Underlying, 4, 27-28, 551\par
\noindent Underlying contract, 439\par
\noindent synthetic options and, 255-260\par
\noindent synthetic short position and, 264-265\par

\marksourcepage{865}
\noindent theoretical value of, 115\par
\noindent volatility independent of price of, 478\par
\noindent Underlying position, 41-43, 103-104\par
\noindent Underlying price\par
\noindent in Black-Scholes model, 66, 345-346\par
\noindent cash index option and, 461-462\par
\noindent exercise price and, 49-51\par
\noindent gamma and changes in, 149-150\par
\noindent lognormal distribution of, 478-481\par
\noindent option value and, 37-38\par
\noindent in theoretical pricing models, 66\par
\noindent theta and changes in, 141-142\par
\noindent vega and changes in, 145-146\par
\noindent Unrealized profits, 7\par
Upside contract position, 419,\par
\noindent Value Line Index, 452\par
\noindent Vanilla option, 551\par
\noindent Vanna, 139, 140, 146, 551\par
\noindent Variance contract, 513\par
\noindent Variance swap, 513\par
\noindent Variation, 8, 551\par
\noindent Vega, 110, 113, 116, 552\par
\noindent in Black-Scholes model, 354-357\par
\noindent of bull spreads, 224\par
\noindent decay, 147, 149, 356-357, 552\par
\noindent at the money, 145-146\par
\noindent negative, 230, 236\par
\noindent neutral, 399\par
\noindent notional, 513\par
\noindent as time passes, 147-148\par
\noindent underlying price changes and, 145-146\par

\marksourcepage{866}
\noindent volatility changes influencing, 146-147\par
\noindent Vega (volatility) risk, 228\par
\noindent Vertical spread, 213-226, 329, 332-333, 552\par
\noindent VIX. See Volatility index\par
\noindent Volatility, 552\par
\noindent in Black-Scholes model, 68, 339\par
\noindent breakeven, 116-117, 130\par
\noindent calculating, 556-557\par
\noindent delta and changes in, 135-136\par
\noindent forecast, 87\par
\noindent forecasting, 391-394\par
\noindent forward, 404-409\par
\noindent future, 394-397\par
\noindent gamma and changes in, 150-151\par
\noindent hedging to reduce, 333-334\par
\noindent implied, 88-95, 512\par
\noindent interpreting data on, 85-95\par
\noindent in-the-money options and, 92\par
\noindent mean reverting, 387-388\par
\noindent observed price changes and, 80-81\par
\noindent realized, 86-87, 116, 512\par
\noindent rising and falling, 468-469\par
\noindent risk, 228-234\par
\noindent scaled for time, 78-79\par
\noindent shifting, 500-501\par
\noindent smirk, 485\par
\noindent as speed of market, 69\par
\noindent as standard deviation, 77-78\par
\noindent surface, 499\par
\noindent term structure of, 390\par
in theoretical pricing models, 80, 89-90, 204, 206, 380\par
\noindent theoretical value and changes in, 146-147\par

\marksourcepage{867}
\noindent theta and changes in, 144-145\par
\noindent underlying contract price and, 478\par
\noindent value, 299, 304\par
\noindent vega and changes in, 146-147\par
\noindent Volatility contracts, 512-514\par
\noindent applications, 536-537\par
\noindent replicating, 534-536\par
\noindent used for hedging, 536\par
\noindent Volatility index (VIX), 515-519\par
\noindent characteristics of, 520-523\par
\noindent futures, 524-529\par
\noindent options, 530-533\par
\noindent trading, 523-524\par
\noindent Volatility skew, 485-486, 532, 552\par
\noindent Volatility smile. See Volatility skew\par
\noindent Volatility spread, 169\par
\noindent adjustments, 206-208\par
\noindent butterfly, 173-176\par
\noindent calendar spread, 182-191\par
\noindent Christmas tree, 182\par
\noindent condor, 176-178\par
\noindent diagonal spread, 196-199\par
\noindent hedgers using, 332\par
\noindent interest rates and dividends in, 192-195\par
\noindent ratio spread, 179-182\par
\noindent straddles, 170-171\par
\noindent strangle, 171-173\par
\noindent submitting an order, 208-209\par
\noindent synthetics used in, 287-290\par
\noindent time butterfly, 191-192\par
\noindent Volga, 147-148, 552\par
\noindent Volume-weighted average price (VWAP), 450\par

\marksourcepage{868}
\noindent Vomma. See Volga\par
\noindent VWAP. See Volume-weighted average price\par
\noindent Warrant, 552\par
\noindent Wilmot, Paul, 57\par
\noindent Wings of butterfly, 192\par
\noindent Write, 325, 552\par
\noindent Yield volatility, 82\par
\noindent Zero-cost collar, 330, 552\par
\noindent Zero-mean assumptions, 382\par
\noindent Zomma, 152-154, 552\par

\marksourcepage{869}
\noindent About the Author\par
Sheldon Natenberg began his trading career in 1982 as an independent\par
market maker in equity options at the Chicago Board Options Exchange. From 1985 to 2000, he traded commodity options, also as an independent floor trader, at the Chicago Board of Trade.\par
While continuing to trade, Mr. Natenberg has also become active as an\par
educator. In this capacity, he has conducted seminars for option traders at major exchanges and professional trading firms in the United States, Europe, and the Far East. In 2000, Mr. Natenberg joined the education team at Chicago Trading Company, a proprietary derivatives trading firm.\par

